\section{Constitutional AI: aprender a partir de AI Feedback}

\subsection{Por qué sustituir una gran cantidad de etiquetas de preferencia con principios}

RLHF requiere una gran cantidad de comparaciones humanas. Constitutional AI intenta comprimir el juicio de valores humano en un pequeño número de principios explícitos, para que el modelo, con base en esos principios, critique, revise y compare respuestas, ampliando así los datos de post-entrenamiento mediante AI feedback.

\videofigure{figures/ai-feedback-scaling.jpg}{Un pequeño número de principios humanos impulsa una gran cantidad de AI critique y preference labels.}{00:46:32--00:48:42}

\begin{knowledgebox}{Una constitution es una norma ejecutable, no una lista de consignas}
Los principios eficaces deben poder aplicarse a una respuesta concreta: señalar qué párrafo viola qué requisito y proponer cómo revisarlo. Los principios demasiado abstractos, que se contradicen entre sí o que carecen de prioridad hacen que el modelo produzca juicios inestables.
\end{knowledgebox}

\subsection{Etapa uno: Critique--Revision Supervised Learning}

El proceso es el siguiente:

\begin{enumerate}
    \item Se usa un red-team prompt para inducir al asistente base a producir una respuesta problemática;
    \item se elige una constitutional principle y se le pide al modelo que critique la respuesta;
    \item se le pide al modelo que reescriba la respuesta con base en la crítica;
    \item se usa la respuesta revisada para hacer supervised fine-tuning.
\end{enumerate}

\videofigure{figures/constitutional-pipeline.jpg}{Constitutional AI usa principios para impulsar la autocrítica y la revisión.}{00:48:42--00:49:38}

\videofigure{figures/constitution-example.jpg}{Ejemplos de critique / revision principle para daño, sesgo y adecuación infantil.}{00:49:38--00:50:48}

\videofigure{figures/constitutional-sft.jpg}{En la primera etapa se usan para SFT las respuestas seguras revisadas por el modelo.}{00:50:48--00:51:24}

\subsection{Beneficios y costos de las revisiones en varias rondas}

A medida que aumenta el número de revisions, la harmlessness mejora, pero la helpfulness puede disminuir: el modelo se vuelve excesivamente reacio a responder por evitar riesgos. El objetivo general no es maximizar la seguridad de forma aislada, sino encontrar una mejora de Pareto entre helpfulness y harmlessness.

\videofigure{figures/revision-results.jpg}{Varias rondas de revision mejoran la harmlessness, pero pueden reducir al mismo tiempo la helpfulness.}{00:51:24--00:52:20}

\begin{warningbox}{Más principios no siempre son mejores}
Varios principios pueden entrar en conflicto entre sí, por ejemplo «intentar responder siempre al usuario» frente a «evitar cualquier riesgo potencial». Si no hay condiciones de aplicación ni prioridades, el modelo optará por la estrategia de rechazo más conservadora, aparentemente segura pero sin utilidad práctica.
\end{warningbox}

\subsection{Etapa dos: RLAIF}

En la etapa de RL, el modelo compara varias respuestas candidatas con base en la constitution y genera AI preference labels; después se entrena un preference model y se optimiza la policy del asistente igual que en RLHF.

\videofigure{figures/rlaif.jpg}{RLAIF usa una AI que, con base en principios, genera preferencias, para después entrenar el reward model y la policy.}{00:52:20--00:56:36}

Si los candidatos son $y_a,y_b$, el AI judge da la siguiente probabilidad de preferencia:

$$
P(y_a\succ y_b\mid x,\mathcal C)
=\sigma\!\left(r_\phi(x,y_a,\mathcal C)-r_\phi(x,y_b,\mathcal C)\right),
$$

\begin{itemize}
    \item $\mathcal C$: constitution;
    \item $r_\phi$: preference model que evalúa con base en los principios;
    \item $\sigma$: función logistic.
\end{itemize}

\subsection{Resultado: llevar los principios a escala como señal de entrenamiento}

\videofigure{figures/constitutional-results.jpg}{Constitutional SFT y RL forman una mejora de Pareto en helpfulness--harmlessness.}{00:56:36--01:00:02}

El valor de Constitutional AI es convertir a los humanos de anotadores muestra por muestra en diseñadores de reglas; pero el AI judge todavía puede malinterpretar los principios, compartir el sesgo del generator o ser sensible a la redacción del prompt.

\begin{importantbox}{El AI feedback debe calibrarse con una evaluación externa}
La autoevaluación del modelo puede ampliar los datos, pero no puede demostrar por sí misma que es correcta. Debe usarse un red-team independiente, muestreo humano, behavior benchmarks y retroalimentación del despliegue real para calibrar continuamente la constitution y el judge, y evitar un ciclo cerrado autoconsistente pero desviado de los objetivos humanos.
\end{importantbox}

\subsection{Resumen de la sección}

Constitutional AI convierte los principios en critique, revision y preference labels, que se usan respectivamente para SFT y RL. Reduce de manera notable el costo de anotación humana muestra por muestra, pero traslada la dificultad al diseño de principios, el manejo de conflictos y la calibración del AI judge.

